\section{卡帕曲线}
\markboth{76. 卡帕曲线}{76. 卡帕曲线}

\noindent
\begin{minipage}[t]{0.50\textwidth}
  \vspace{0pt}
  \hspace*{2em}如图 76.1 所示的曲线，形状很像一个希腊字母 $\kappa$（读作“卡帕”），所以叫做卡帕曲线。它的极坐标方程是
  \[
    \rho = a\ctg\theta.
  \]
  \hspace*{2em}这种曲线可用下面的几何方式定义：设动圆 $C$ 的半径等于定长 $a$，圆心 $C$ 在 $x$ 轴上滑动，那么从坐标原点 $O$ 向圆 $C$ 所作切线的切点 $M$ 的轨迹叫做卡帕曲线。
\end{minipage}
\hfill
\begin{minipage}[t]{0.47\textwidth}
  \vspace{2pt}
  \centering
  \begin{tikzpicture}[scale=0.82, >=Stealth, line cap=round, line join=round]
    \def\a{0.92}
    \def\th{35}
    \pgfmathsetmacro{\xc}{\a/sin(\th)}
    \pgfmathsetmacro{\xm}{\a*cos(\th)/sin(\th)*cos(\th)}
    \pgfmathsetmacro{\ym}{\a*cos(\th)}

    % 坐标轴 (右端足够长，避开虚线圆)
    \draw[->] (-2.15, 0) -- (2.82, 0);
    \node[below=1.5pt, font=\scriptsize] at (2.74, 0) {$x$};
    \draw[->] (0, -1.25) -- (0, 1.38) node[left, font=\scriptsize] {$y$};

    % 渐近线 y = +/- a
    \draw[densely dashed] (-2.08, \a) -- (2.48, \a);
    \draw[densely dashed] (-2.08, -\a) -- (2.48, -\a);

    % 动圆 C (虚线)
    \draw[densely dashed] (\xc, 0) circle (\a);

    % 卡帕曲线四条分支 (左右上下对称延伸至渐近线附近)
    \draw[very thick, curveblue, domain=22:158, samples=160, variable=\t]
      plot ({\a*cos(\t)*cos(\t)/sin(\t)}, {\a*cos(\t)});
    \draw[very thick, curveblue, domain=202:338, samples=160, variable=\t]
      plot ({\a*cos(\t)*cos(\t)/sin(\t)}, {\a*cos(\t)});

    % 切线 OM 与半径 CM，切点直角记号
    \coordinate (O) at (0,0);
    \coordinate (C) at (\xc, 0);
    \coordinate (M) at (\xm, \ym);
    \draw (O) -- ($(O)!1.25!(M)$);
    \draw (C) -- (M) node[midway, above right=-1pt, font=\scriptsize] {$a$};
    \draw ($(M)+(\th+180:0.14)$) -- ++(\th-90:0.14) -- ($(M)+(\th-90:0.14)$);

    % 标注角 theta
    \draw (0.36, 0) arc (0:\th:0.36);
    \node[font=\scriptsize] at (17:0.50) {$\theta$};

    \node[below left=-1pt, font=\scriptsize] at (O) {$O$};
    \fill (C) circle (1.2pt) node[below=1pt, font=\scriptsize] {$C$};
    \fill (M) circle (1.2pt) node[above=2.5pt, font=\scriptsize] {$M$};
  \end{tikzpicture}
  \captionof*{figure}{图\ 76.1}
\end{minipage}

把卡帕曲线极坐标方程的两边同乘以 $\rho\sin\theta$，然后平方，得到
\[
  \rho^2\cdot(\rho\sin\theta)^2 = a^2(\rho\cos\theta)^2.
\]
化成直角坐标，得到
\[
  y^2(x^2+y^2) = a^2x^2,
\]
这就是卡帕曲线的直角坐标方程。

卡帕曲线是十七世纪发现的，在史留斯和惠更斯的通信中第一次提到这种曲线。当时是为了解答下面的运动学问题：同一平面内有一个固定的直角坐标系 $Oxy$ 和一个运动的直角坐标系 $O'x'y'$，已知 $x'$ 轴始终通过点 $O$，$y'$ 轴与 $x$ 轴的交点 $C$ 到 $O'$ 的距离等于定长 $a$，求动坐标系的原点 $O'$ 的轨迹。容易看出，这里的点 $O'$ 相当于图 76.1 中的 $M$，$x'$ 轴相当于直线 $MO$，因而所求的轨迹就是图 76.1 中的卡帕曲线。
